\begin{solapartat}
\punts{0,25} $\Delta s = \dfrac{1}{2}gt^2 \Longrightarrow g = \dfrac{2\Delta s}{t^2} = \dfrac{2\cdot 1}{1,1^2} = 1,65\ \mathrm{m/s^2}$.

\punts{0,75}

$\displaystyle g = G\,\frac{M_{Lluna}}{R_{Lluna}^2} \Rightarrow M_{Lluna} = g\,\frac{R_{Lluna}^2}{G} = 1,65\,\frac{(1,74\times 10^{6})^2}{6,67\times 10^{-11}} = 7,50\times 10^{22}\ \mathrm{kg}$ \quad \punts{0,25}
\end{solapartat}

\begin{solapartat}
La segona llei de Newton estableix que: $\vec{F} = m\vec{a}$ \quad \punts{0,1}

Segons la llei de gravitació universal, el mòdul de la força s'expressa com:

$F = G\,\dfrac{mM_T}{r^2}$ \quad \punts{0,4}

Per tant, obtenim que: $a = G\,\dfrac{M_T}{r^2}$ \quad \punts{0,2}

D'altra banda, considerant que la Lluna descriu un moviment circular uniforme al voltant de la terra, la seva acceleració centrípeta és: $a = \dfrac{v^2}{r}$ o $a = \omega^2 r$ \punts{0,25}\unskip, i la velocitat es pot expressar com $v = \dfrac{2\pi r}{T}$ o $\omega = \dfrac{2\pi}{T}$ \quad \punts{0,1}

Per tant: $T = 2\pi\sqrt{\dfrac{r^3}{GM_T}}$ \quad \punts{0,1}

Llavors: $T = 2\pi\sqrt{\dfrac{(3,84\times 10^{8})^3}{6,67\times 10^{-11}\cdot 5,98\times 10^{24}}} = 2,37\times 10^{6}\ \mathrm{s} = 27,4\ \mathrm{dies}$ \quad \punts{0,1}
\end{solapartat}
